% !PRINTED 314
\weberpage{314}
%<<CONT
further the quantity $A$ defined by (1) of the previous section; but then also the symmetric functions of the seven quantities $\Psi^2$, and further, if we now add for the sake of simplicity the assumption, to which we shall return, that none of the quantities $\Psi$ vanishes, the symmetric functions of the quantities $v$; and not only the symmetric functions of the $v$, but all functions of the $v$ that admit the permutations of the group $T$.

The seven quantities

\weq{2}{\begin{aligned}
v_{1,0,0} &= v_1, & v_{0,1,0} &= v_2, & v_{0,0,1} &= v_3, \\
v_{0,1,1} &= v_4, & v_{1,0,1} &= v_5, & v_{1,1,0} &= v_6, & v_{1,1,1} &= v_7
\end{aligned}}

are then the roots of a rational equation of degree $7$ with the group $T$, and by addition of equation (1) and the equations [\S.\,80, (1)]:

\[
\Sigma (-1)^{\Sigma z\xi} X_{z_1, z_2, z_3} = \Psi_{\xi_1, \xi_2, \xi_3}
\]

there results, having regard to (9) of the previous section:

\weq{3}{8\,X_{0,0,0} = A\Bigl\{\begin{aligned}[t]
&\sqrt{v_1}\sqrt{v_5}\sqrt{v_6}\sqrt{v_7} + \sqrt{v_2}\sqrt{v_4}\sqrt{v_6}\sqrt{v_7} \\
&+ \sqrt{v_3}\sqrt{v_4}\sqrt{v_5}\sqrt{v_7} + \sqrt{v_2}\sqrt{v_3}\sqrt{v_5}\sqrt{v_6} \\
&+ \sqrt{v_1}\sqrt{v_3}\sqrt{v_4}\sqrt{v_6} + \sqrt{v_1}\sqrt{v_2}\sqrt{v_4}\sqrt{v_5} \\
&+ \sqrt{v_1}\sqrt{v_2}\sqrt{v_3}\sqrt{v_7}
\end{aligned}\Bigr\} + B.}

There are 15 changes of sign of the seven square roots $\sqrt{v}$ under which this expression remains unchanged, namely when all seven signs are changed simultaneously, or when the signs of a triple or of a quadruple are changed. \gesperrt{Consequently the expression (3) takes only eight different values, however the seven square roots occurring in it may be determined.}

\begin{center}
\S.\,83.

\gesperrt{Metacyclic Equations of Degree $8$.}
\end{center}

Now in order, after these preparations, to find the primitive metacyclic equations of degree $8$, we have first to determine the metacyclic divisors of the group $S$. The group $S$ itself is, as already mentioned, not metacyclic. Its
%CONT>>
